\chapter*{Conclusions}
\addcontentsline{toc}{chapter}{Conclusions}
\markboth{Conclusions}{Conclusions}

This thesis set out to examine a simple but demanding idea: that a vehicle can decide from one camera
image, without a map, whether each of its possible maneuvers at the next intersection is permitted by
the traffic signals. The main objective was to establish, independently and from public data only, how
well current models solve this task, what their accuracy depends on, and how far it carries beyond the
data they were trained on.

\paragraph{Answers to the research questions.}
\begin{description}[leftmargin=1.2cm,style=nextline]
  \item[RQ1] \emph{Benchmark reconstruction.} Yes. Labels rebuilt from the OpenLane-V2 lane and
        traffic-element topology reproduce the published per-direction label distribution for almost
        every class, with one systematic difference (left \emph{none}). The evidence also indicates that
        the original benchmark combines subsets A and B under US rules, rather than subset A alone as
        stated.
  \item[RQ2] \emph{Structured model.} A structured pipeline---an RF-DETR detector whose classes encode
        the maneuvers each light controls, followed by explicit rules---reaches a TSS of 0.899 on subset A
        and 0.915 on subset B, in the range published for the end-to-end direction-aware model, although
        the two are not directly comparable. It fails where the data are thin: right-turn reds, yellow,
        and straight maneuvers that no light governs.
  \item[RQ3] \emph{Transfer.} A model trained on German data (DTLD, test TSS 0.958) transfers well to US
        day-time driving (LISA day 0.952 without right turns) but loses ten points at night and doubles
        its False-Go rate on Argoverse~2 images. Detection of small lights is the dominant failure: on
        ATLAS, recall falls from over 85\% for lights taller than 40 pixels to about 20\% for lights of
        10--20 pixels.
\end{description}

\paragraph{What emerged from the analysis.} Beyond the individual numbers, the work points to three
broader lessons. First, aggregate scores conceal the decisions that matter most: models with
state-of-the-art TSS miss almost every red right-turn arrow, and yellow lights and straight maneuvers
without a governing signal are weak points as well. Second, traffic rules are not a detail of the
labels but part of the task's definition; the same image calls for different decisions in the United
States and in Germany, and a design that keeps the rules explicit can adapt to this, where an end-to-end
model has to be relabeled and retrained. Third, in the structured pipeline an undetected light silently becomes a \emph{go}
decision, so robustness out of distribution is above all a question of detecting small, distant lights
and of treating their absence as uncertainty rather than permission.

\paragraph{Limitations.} The conclusions rest on reconstructed labels without full manual quality
control, on single training runs per configuration, on evaluation with the validation split of
OpenLane-V2, and on single-frame models. Small differences between models are therefore reported as
non-significant rather than as rankings.

\paragraph{Future developments.} The most promising continuations are safety-weighted training and
evaluation protocols that report per-maneuver False-Go rates, models conditioned on the applicable
traffic convention, temporal aggregation over short sequences, training on mixed datasets from several
countries and lighting conditions, and better detection of small lights through resolution or a tele
camera. Together they would move maneuver-level signal understanding from a benchmark score towards a
component whose failures are understood, bounded and acceptable on the road.
